% !TeX root = 04_intro_ml_polytech.tex
% !TeX encoding = UTF-8
% !TeX spellcheck = en_GB
\begin{frame}{Neural Network Types}{What kind of inductive bias?}
    We talked about the \textbf{bias} of a model:

    \vfill

    \begin{center}
        \def\overfitcurve{(0.30,1.10) (0.37,1.97) (0.45,1.90) (0.52,1.54) (0.59,1.25) (0.67,1.15) (0.74,1.25) (0.81,1.48) (0.89,1.77) (0.96,2.05) (1.03,2.25) (1.11,2.37) (1.18,2.38) (1.25,2.31) (1.33,2.18) (1.40,2.02) (1.47,1.86) (1.55,1.74) (1.62,1.67) (1.69,1.65) (1.77,1.69) (1.84,1.77) (1.92,1.86) (1.99,1.96) (2.06,2.04) (2.14,2.09) (2.21,2.09) (2.28,2.05) (2.36,1.98) (2.43,1.88) (2.50,1.78) (2.58,1.69) (2.65,1.62) (2.72,1.57) (2.80,1.52) (2.87,1.44) (2.94,1.33) (3.02,1.14) (3.09,0.92) (3.16,0.77) (3.20,0.80)}
\def\goodcurve{(0.30,1.05) (0.50,1.31) (0.70,1.54) (0.90,1.72) (1.10,1.87) (1.30,1.98) (1.50,2.04) (1.70,2.06) (1.90,2.05) (2.10,1.99) (2.30,1.88) (2.50,1.74) (2.70,1.55) (2.90,1.31) (3.10,1.03) (3.20,0.88)}
\def\underfitcurve{(0.30,1.69) (3.20,1.54)}
\def\trainpts{0.30/1.10, 0.62/1.19, 0.94/1.99, 1.27/2.29, 1.59/1.70,
    1.91/1.86, 2.23/2.08, 2.56/1.71, 2.88/1.43, 3.20/0.80}
\def\testpts{0.36/0.99, 0.98/1.68, 1.19/2.13, 2.10/2.21,
    2.52/1.85, 2.65/1.47, 2.80/1.53, 2.86/1.63}

\begin{tikzpicture}[
        train/.style={fill=CEAred, draw=none, circle, inner sep=1.6pt},
        test/.style={fill=CEAblue, draw=none, circle, inner sep=1.6pt},
        ax/.style={thick, CEAblack, ->}
    ]
    \foreach \xs/\firsttitle/\subtitle/\title/\curve/\col/\msetrain/\msetest/\name/\effect/\adjective in {
            0/Model 1/Overfitting/High Variance/\overfitcurve/CEAblack/0.0/0.153/overfitting/memorised the training data/unreliable,
            4.7/Model 2/Generalisation/Optimal/\goodcurve/CEAblack/0.044/0.030/generalisation/adjusted to the training data/great,
            9.4/Model 3/Underfitting/High Bias/\underfitcurve/CEAblack/0.200/0.148/underfitting/not learnt the training data/useless}
        {
            \begin{scope}[xshift=\xs cm]
                \draw[ax] (-0.2,0) -- (3.6,0) node[below] {$x$};
                \draw[ax] (0,-0.2) -- (0,3.0) node[left] {$y$};
                \node<1-2>[above] at (1.7,2.9) {\textbf{\firsttitle}};
                \node<3->[above] at (1.7,2.9) {\textcolor{CEAred}{\textbf{\title}}};
                \node<3->[above] at (1.7,3.4) {\footnotesize\textcolor{CEAred}{\textbf{\subtitle}}};
                \draw[\col, thick, smooth] plot coordinates {\curve};
                \foreach \px/\py in \trainpts \node[train] at (\px,\py) {};
                \only<2->{%
                    \foreach \px/\py in \testpts \node[test] at (\px,\py) {};
                }
            \end{scope}
        }
\end{tikzpicture}

    \end{center}

    \vfill

    And we saw that \textbf{higher bias} leads to a more constrained model, and that sometimes it can be more adapted to describe the data (e.g.\ linear data)
\end{frame}

\begin{frame}{Neural Network Types}{What kind of inductive bias?}
    How can we impose some \textbf{inductive bias} on our model? What kind of operations are appropriate for different data?

    \vfill

    \begin{columns}[T]
        \begin{column}{0.5\textwidth}
            \centering
            \def\curve{(0.30,1.10) (0.37,1.97) (0.45,1.90) (0.52,1.54) (0.59,1.25) (0.67,1.15) (0.74,1.25) (0.81,1.48) (0.89,1.77) (0.96,2.05) (1.03,2.25) (1.11,2.37) (1.18,2.38) (1.25,2.31) (1.33,2.18) (1.40,2.02) (1.47,1.86) (1.55,1.74) (1.62,1.67) (1.69,1.65) (1.77,1.69) (1.84,1.77) (1.92,1.86) (1.99,1.96) (2.06,2.04) (2.14,2.09) (2.21,2.09) (2.28,2.05) (2.36,1.98) (2.43,1.88) (2.50,1.78) (2.58,1.69) (2.65,1.62) (2.72,1.57) (2.80,1.52) (2.87,1.44) (2.94,1.33) (3.02,1.14) (3.09,0.92) (3.16,0.77) (3.20,0.80)}
\begin{tikzpicture}
    % Tabular data
    \draw<2>[thick, CEAred, fill=CEAred!15, xstep=1, ystep=0.5] (0, 3.5) rectangle (6, 4);
    \draw<2>[thick, CEAblack, fill=none, xstep=1, ystep=0.5] (0, 0) grid (6, 4);
    \draw<2>[thick, <->, CEAblue] (0.5, 1.75) -- (2.5, 1.75);
    \draw<2>[thick, <->, CEAblue] (0.5, 2.75) -- (3.5, 2.75);
    \draw<2>[thick, <->, CEAblue] (1.5, 0.75) -- (4.5, 0.75);

    % Image/Video
    \draw<3>[thick, CEAblack, xstep=0.2, ystep=0.2] (0, 0) grid (3, 3);
    \draw<3>[thick, CEAblue, fill=CEAblue!25] (0, 3) -- (3, 3) -- (3, 0) -- (3.33, 0.33) -- (3.33, 3.33) -- (0.33, 3.33) -- cycle;
    \draw<3>[thick, CEAgreen, fill=CEAgreen!25] (0.33, 3.33) -- (3.33, 3.33) -- (3.33, 0.33) -- (3.66, 0.66) -- (3.66, 3.66) -- (0.66, 3.66) -- cycle;
    \draw<3>[thick, CEAred, fill=CEAred!25] (0.66, 3.66) -- (3.66, 3.66) -- (3.66, 0.66) -- (4, 1) -- (4, 4) -- (1, 4) -- cycle;
    \draw<3>[thick, CEAblack] (0, 3) -- (1, 4) -- (4, 4) -- (3, 3);
    \draw<3>[thick, CEAblack] (3, 0) -- (4, 1) -- (4, 4);
    \draw<3>[thick, CEAyellow, fill=CEAyellow, fill opacity=0.35] (1, 1) circle (0.5);
    \draw<3>[thick, CEAyellow, fill=CEAyellow, fill opacity=0.35] (2, 1.5) circle (0.5);
    \draw<3>[thick, CEAyellow, fill=CEAyellow, fill opacity=0.35] (0.75, 2.5) circle (0.5);

    % Sequences
    \draw<4->[thick, ->, CEAblack] (-0.1, 0.0) -- (3.0, 0.0) node[below] {$\lambda$};
    \draw<4->[thick, ->, CEAblack] (0.0, -0.1) -- (0.0, 3.0) node[left] {$S$};
    \draw<4->[CEAblack, thick, smooth] plot coordinates {\curve};

    \node<4->[anchor=north, align=center, CEAblack] at (1.5, -1.0) {The cat is on the table};
\end{tikzpicture}

        \end{column}
        \hfil
        \begin{column}{0.5\textwidth}
            \onslide<2->{%
                \textbf{Tabular data:}
                \begin{itemize}
                    \item Every column may influence other columns
                    \item No concept of neighbourhood of columns
                \end{itemize}
            }
            \onslide<3->{%
                \textbf{Images/Videos:}
                \begin{itemize}
                    \item Information is locally connected
                    \item The same local operation applies everywhere (\emph{weight sharing})
                \end{itemize}
            }
            \onslide<4->{%
                \textbf{Sequences:}
                \begin{itemize}
                    \item Information depends on distance
                    \item Next bit depends on previous
                \end{itemize}
            }
        \end{column}
    \end{columns}
\end{frame}

\begin{frame}{Neural Network Types}{Convolution operations}
    Other than just \textbf{matrix multiplication} in fully connected networks
    \begin{equation*}
        \hat{y} =
        f^{(L)}
        \left(
        f^{(L-1)}
        \left(
            \dots
            W^{(L-1)}
            \right)
        W^{(L)}
        \right)
    \end{equation*}
    we can also have \emph{sliding doors} operations:

    \begin{center}
        \begin{tikzpicture}
    \foreach \x [count=\i from 0] in {0, 4cm, 8cm} {%
            \begin{scope}[xshift=\x]
                \draw[CEAred, fill=CEAred!25] (0+\i*0.5, 0) rectangle (1.5+\i*0.5, 1);
                \draw[CEAred] (0+\i*0.5, 0) -- (0.5+\i*0.5, -0.5);
                \draw[CEAred] (1.0+\i*0.5, -0.5) -- (1.5+\i*0.5, 0.0);
                \draw[CEAgreen, fill=CEAgreen!15, thick, xstep=0.5] (0, 0) grid (3, 1);
                \draw[CEAgreen, fill=CEAgreen!15, fill opacity=0.25, thick] (0.5, -1.5) rectangle (2.0, -0.5);
                \draw[CEAgreen, fill=CEAgreen!15, fill opacity=0.25, thick] (1.0, -1.5) -- (1.0, -0.5);
                \draw[CEAgreen, fill=CEAgreen!15, fill opacity=0.25, thick] (1.5, -1.5) -- (1.5, -0.5);
                \draw[CEAgreen, fill=CEAgreen!25, thick] (0.5+\i*0.5, -1.5) rectangle (1.0+\i*0.5, -0.5);
            \end{scope}
        }
\end{tikzpicture}

    \end{center}

    \vfill

    They enable to \textbf{aggregate} local information to reduce information.
\end{frame}

\begin{frame}{Neural Network Types}{Convolution operations}
    Formally, they are defined as:

    \begin{equation*}
        \text{(1D)}~
        \tilde{S}^{(c)}_{i}
        =
        \sum\limits_{c^{\prime}}
        \sum\limits_{u = -\lfloor K_u/2 \rfloor}^{\lfloor K_u/2 \rfloor}
        \underbrace{\mathcal{K}^{(c c^{\prime})}_{u}}_{\text{kernel}}\;
        S^{(c^{\prime})}_{i-u}
        \qor \text{(2D)}~
        \tilde{X}^{(c)}_{ij}
        =
        \sum\limits_{c^{\prime}}
        \sum\limits_{u = -\lfloor K_u/2 \rfloor}^{\lfloor K_u/2 \rfloor}
        \sum\limits_{v = -\lfloor K_v/2 \rfloor}^{\lfloor K_v/2 \rfloor}
        \underbrace{\mathcal{K}^{(c c^{\prime})}_{uv}}_{\text{kernel}}\;
        X^{(c^{\prime})}_{i-u,j-v}
    \end{equation*}

    and they enable different operations on local clusters of \emph{pixels}:

    \begin{center}
        \begin{tikzpicture}[scale=0.8, every node/.style={scale=0.8}]
    \draw[thick, CEAblack, xstep=0.2, ystep=0.2] (0, 0) grid (3, 3);
    \draw[thick, CEAblue, fill=CEAblue!25] (0, 3) -- (3, 3) -- (3, 0) -- (3.33, 0.33) -- (3.33, 3.33) -- (0.33, 3.33) -- cycle;
    \draw[thick, CEAgreen, fill=CEAgreen!25] (0.33, 3.33) -- (3.33, 3.33) -- (3.33, 0.33) -- (3.66, 0.66) -- (3.66, 3.66) -- (0.66, 3.66) -- cycle;
    \draw[thick, CEAred, fill=CEAred!25] (0.66, 3.66) -- (3.66, 3.66) -- (3.66, 0.66) -- (4, 1) -- (4, 4) -- (1, 4) -- cycle;
    \draw[thick, CEAblack] (0, 3) -- (1, 4) -- (4, 4) -- (3, 3);
    \draw[thick, CEAblack] (3, 0) -- (4, 1) -- (4, 4);
    \draw[very thick, CEAred, fill=CEAred, fill opacity=0.5] (1, 1) rectangle (2.0, 2.0);
    \draw[ultra thick, CEAred, ->] (2.2, 1.5) -- (3.2, 1.5);
    \draw[ultra thick, CEAred, ->] (0.8, 1.5) -- (-0.2, 1.5);
    \draw[ultra thick, CEAred, ->] (1.5, 2.2) -- (1.5, 3.2);
    \draw[ultra thick, CEAred, ->] (1.5, 0.8) -- (1.5, -0.2);
\end{tikzpicture}

    \end{center}
\end{frame}

\begin{frame}{Neural Network Types}{Convolutional Neural Networks}
    The operation is usually defined as $X_{ij} = (\mathcal{K} * Z)_{ij}$ and a composition of such operations and \textbf{activation functions}:

    \begin{equation*}
        \hat{y}
        =
        f^{(L)}
        \left(
        \mathcal{K}^{L} *
        f^{(L-1)}
        \left(
            \mathcal{K}^{(L-1)} *
            \dots
            \right)
        \right)
    \end{equation*}

    is called (fully) \textbf{convolutional neural network} (CNN).


    \vfill \pause

    Fully connected networks and CNNs are often coupled to create powerful models:

    \begin{center}
        \includegraphics[width=0.6\textwidth]{04_slides/img/AlexNet_architecture.png}
    \end{center}
\end{frame}

\begin{frame}{Neural Network Types}{\enquote{Temporal} sequence modelling}
    Modern AI deals with a lot of \textbf{sequence} modelling (e.g.\ LLMs for text, time series, etc.).
    Two main ingredients are necessary:
    \begin{itemize}
        \item Introduce the concept of \enquote{position} in a \enquote{sentence}
        \item Study how related are two concepts in the same \enquote{sentence}
    \end{itemize}

    \vfill \pause

    \begin{center}
        \begin{tikzpicture}
    \draw[CEAblack, fill=CEAblack!15] (0, 0) rectangle (1, 4) node[midway, rotate=90] (data) {DATA};

    \draw[CEAblack, ->, ultra thick] (1.1, 2) -- (1.9, 2);

    \draw[CEAblue, fill=CEAblue!15] (2.1, 0) rectangle (3.1, 4) node[midway, rotate=90] (pe) {\small POSITIONAL EMBEDDING};

    \draw[CEAblack, ->, ultra thick] (3.2, 2) -- (4.1, 2);

    \draw[CEAred, fill=CEAred!15] (4.2, 0) rectangle (5.2, 4) node[midway, rotate=90] (att) {\small ATTENTION MECHANISM};

    \node[anchor=west] at (5.3, 2) {\dots};
\end{tikzpicture}

    \end{center}
\end{frame}

\begin{frame}{Neural Network Types}{\enquote{Temporal} sequence modelling}
    The mechanism uses the same idea that you use to find important documents/notes:

    \vfill

    \begin{enumerate}
        \item<1-> Formulate a sentence/\textbf{sequence} and assign a \enquote{reading direction} (e.g.\ left to right)
        \item<2-> Identify important parts of the sentence:
            \begin{itemize}
                \item \textbf{Query} (Q): the information you want to retrieve
                \item \textbf{Key} (K): how the information can be looked for in Q
                \item \textbf{Value} (V): the information content once ranked by Q and K
            \end{itemize}
        \item<3-> \textbf{cross-correlate Q and K}: find how the K matches your request Q
        \item<4-> \textbf{find the right direction V}: project your search in the direction of \enquote{value}
    \end{enumerate}

    \vfill

    \begin{onlyenv}<1>
        \begin{block}{Input data}
            \begin{equation*}
                X \gets X + \operatorname{PositionalEncoding}(1, \dots, S) \in \mathbb{R}^{S \times D}
            \end{equation*}
        \end{block}

        {\tiny $X$ is multi-dimensional in general. If you think of LLMs, imagine each word as an integer $Z \in \mathbb{N}^{S}$, which is then assigned to a vector of dimension $D$}
    \end{onlyenv}

    \begin{onlyenv}<2>
        \begin{block}{Query, Key, Value}
            \begin{equation*}
                Q = \operatorname{Linear}_Q(X) \in \mathbb{R}^{S \times D_{\text{emb}}},
                \quad
                K = \operatorname{Linear}_K(X) \in \mathbb{R}^{S \times D_{\text{emb}}},
                \quad
                V = \operatorname{Linear}_V(X) \in \mathbb{R}^{S \times D_{\text{emb}}}
            \end{equation*}
        \end{block}
    \end{onlyenv}

    \begin{onlyenv}<3>
        \begin{block}{Attention Matrix}
            \begin{equation*}
                A = \operatorname{softmax}\left( \frac{Q\, K^{\top}}{\sqrt{D_{\text{emb}}}} \right) \in \mathbb{R}^{S \times S}
            \end{equation*}
        \end{block}
    \end{onlyenv}

    \begin{onlyenv}<4>
        \begin{block}{Value Matrix and Output}
            \begin{equation*}
                X^{\prime} = \operatorname{FFN}(A\, V) \in \mathbb{R}^{S \times D_{\text{emb}}}
            \end{equation*}
            where $\operatorname{FFN}$ is a fully-connected (or \emph{feed forward}) neural network.
        \end{block}
    \end{onlyenv}
\end{frame}

\begin{frame}{Neural Network Types}{MultiHead Attention and Transformers}
    This mechanism is known as \textbf{attention mechanism} and can be done multiple times in parallel for each $X$ (multiple \enquote{attention heads} can be \textbf{concatenated} to form a \enquote{multi-attention head}).

    \vfill

    This mechanism is the basic block of \textbf{transformers} used for LLMs, time series modelling, etc.:

    \begin{center}
        \includegraphics[width=0.3\textwidth, angle=90]{04_slides/img/transformer.jpg}
    \end{center}
\end{frame}
